
This work connects three research threads: multi-dimensional trustworthiness evaluation, calibration and uncertainty quantification, and latent factor analysis in machine learning.

\subsubsection{Multi-Dimensional Trustworthiness Evaluation}

Recent frameworks conceptualize LLM trustworthiness as multi-faceted. Zhou et al. (2024) propose the Trust-RAG Compass with six dimensions: factuality, robustness, fairness, transparency, accountability, and privacy. This framework provides conceptual clarity but does not investigate cross-dimensional relationships. TrustLLM \citep{sun2024trustllm} evaluates models across trust dimensions without analyzing whether performance on one predicts performance on others.

Benchmark suites including HELM \citep{liang2023helm} and the Open LLM Leaderboard \citep{beeching2023openllm} report multi-benchmark scores while implicitly treating dimensions as independent. Analysis of 4,561 models reveals cross-benchmark correlations of $\rho$ = 0.80--0.87, suggesting this independence assumption warrants examination.

\subsubsection{Calibration and Representation Stability}

Khanmohammadi et al. (2025) introduce CCPS (Calibrating via Perturbed Stability), achieving 55\% ECE reduction by leveraging internal representation stability. Their finding---that perturbation-stable representations yield better calibration---informs the mechanism hypothesis explored here. However, CCPS examines single-dimension calibration without cross-benchmark analysis.

\subsubsection{Latent Factor Analysis in ML Evaluation}

In psychometrics, Spearman's g-factor explains positive correlations across cognitive tests. Schumacher et al. (2024) analyze benchmark correlation structure and find that model scale explains substantial shared variance. The present work extends their analysis by residualizing on scale and training recency before extracting latent factors, examining whether structure persists beyond what scale alone explains.
